% =====================================================================
%  Harald Høffding, NEW INVESTIGATIONS INTO THE PRIMITIVE FAMILY [Nye
%  Undersøgelser om den primitive Familje], review of C. N. Starcke, Die
%  primitive Familie in ihrer Entstehung und Entwickelung (Leipzig:
%  Brockhaus, 1888).
%  Tilskueren. Maanedsskrift, 5. Aargang (1888), pp. 787--795.
%  TRANSLATION (English)
%
%  Translated from transcription.tex (Danish) in the parent directory,
%  whose header records how its text was established from the Internet
%  Archive scan (every page transcribed at the image, diffed against an
%  independent OCR witness, and read a second time at the image).  No
%  published translation was consulted (none is known).
%
%  ---- REGISTER ----
%  Scholarly, moderately literal, readable English; American spelling.
%  Long periods are broken only where English will not carry them, never
%  across a page marker.
%
%  ---- TERMINOLOGY ----
%  Fixed by the brief: Familje = family; Ægteskab = marriage; Monogami,
%    Polygami, Exogami = monogamy, polygamy, exogamy; Promiskuitet =
%    promiscuity; Slægtskab = kinship; Slægtskabsbestemmelse =
%    determination of kinship; Klan / Stamme = clan / tribe; Moder /
%    Fader = mother / father; Faderforhold = paternal relation;
%    patriarkalsk = patriarchal; Elskovsfølelse = the feeling of erotic
%    love; Kærlighed = love; Sympati = sympathy; Skik = custom; Sæder =
%    mores; sædelig = moral; Etik, etisk = ethics, ethical; Vurdering =
%    valuation; Udvikling = development; Udviklingsproces = process of
%    development; Institution = institution; Motiv = motive;
%    Menneskeslægten = the human race; Folkeslag = peoples; Etnolog =
%    ethnologist.
%  From the repo's translation of Positivisme og Materialisme (1871):
%    Aand, aandelig = spirit, spiritual; Forestilling = representation
%    (so Synsforestillinger = visual representations, Rumsforestilling =
%    spatial representation, Forestillingslivet = the life of
%    representation).
%  Added here: Kvindelinje / Mandslinje = the female line / the male line;
%    de Vilde, vilde Folkeslag = the savages, savage peoples; Forsker =
%    investigator (Undersøgelse = investigation); Kønsinstinkt = the sexual
%    instinct; kønslig = sexual; Parring i Flæng = indiscriminate mating;
%    Udviklingstrin = stage of development; Samfund = society; Løshed /
%    Ubundethed = looseness / unrestraint; Retsforhold / Blodsforhold =
%    legal relation / blood relation; Faderfølelse = paternal feeling;
%    Avlingsforhold = relation of procreation; Kyskhedsfordring = demand
%    of chastity; Moment = factor; Slægt = kin (p. 791), the race (p.
%    794); Elskovssang, Elskovsdrik = love song, love potion.
%
%  ---- CONVENTIONS ----
%  - Page markers \opage{N} and "% --- p. N ---" at the same joints as in
%    transcription.tex; where the page turns inside a Danish word or
%    phrase the English turns at the nearest joint (% TR: at the site).
%  - The same head, bibliographic line (kept exactly as printed, German
%    title, imprint and „66. Bind“), rules, paragraphs, \emph spans
%    (letterspacing in the original, carried over the English words that
%    render the spaced Danish), signature (small caps, set right) and
%    closing double rule as the Danish.
%  - Quotation marks »...« -> ``...''.  German titles and the Latin maxim
%    (p. 788) stay as printed, in roman as the Danish sets them.  Andrew
%    Lang (p. 793) and Starcke (pp. 794--795), whom Høffding quotes in
%    Danish, are translated from his Danish; the English and German
%    originals are not reconstructed.
%  - Translator's decisions are recorded in "% TR:" comments at the site.
%
%  ---- DEFECTS: DANISH AS PRINTED, ENGLISH CORRECT ----
%  The printer's errors listed in the transcription's header are rendered
%  in their evident sense and logged in % TR: at each site: p. 788
%  „nøgterne“ (for nøgternt); p. 789 „Untermehung“ in Bachofen's German
%  title (kept as printed, as is the title's „religiøsen“) and
%  „Promiskniteten“; p. 792 „Kønsinktinktet“; p. 795 the full stop before
%  lower-case „skilles“.  Also p. 788 „Henstillinger“, which the
%  transcription notes as not making sense (rendered as Forestillinger).
%
%  ---- HOW THE TRANSLATION WAS CHECKED (2026-09-27) ----
%  (1) The translator re-read the English against the Danish sentence by
%  sentence for omissions, inversions of sense, dropped negations and
%  modals, and lost concessives (vel, dog, jo, nok, ganske vist).
%  (2) Mechanical audit (tr_audit.py): per printed page the same \opage
%  joints, \emph and \textit spans, centre blocks and paragraph breaks as
%  transcription.tex, and an English/Danish word ratio within 0.95--1.60.
%  (3) Independent review: a reader who translated none of it read every
%  sentence of the Danish against the English and proposed 6 corrections,
%  all applied.  Sense: p. 789 „vist nok“ = probably (Høffding's guess at
%  the investigators' motives, not a certainty); p. 791 „maa her vel
%  fastholdes“ = must presumably be retained (the modal restored); p. 787
%  „ikke altid sagt“ = not always a given (not "by no means").  English:
%  pp. 787, 789, 793.  The reviewer agreed with the p. 788 reading
%  Forestillinger.
% =====================================================================
\documentclass[12pt,a4paper,openany]{book}
\usepackage[utf8]{inputenc}
\usepackage[T1]{fontenc}
\usepackage{textalpha}
\usepackage{libertinus}
\usepackage{libertinust1math}
\usepackage{graphicx}
\usepackage[english]{babel}
\usepackage[protrusion=true,expansion=true]{microtype}
\usepackage{setspace}
\usepackage[margin=1.2in]{geometry}
\usepackage{fancyhdr}
\setlength{\headheight}{28pt}
\addtolength{\topmargin}{-13.5pt}

% Footnotes as the printing sets them: *).
\renewcommand{\thefootnote}{*)}

% The article's numbered sections: the number alone, centred on its line.
\newcommand{\secnum}[1]{\par\begin{center}#1\end{center}\par\nobreak}

\usepackage{hyperref}
\hypersetup{
  pdftitle={New Investigations into the Primitive Family},
  pdfauthor={Harald Høffding},
  hidelinks
}

\pagestyle{fancy}
\fancyhf{}
\fancyhead[LE,RO]{\thepage}
\fancyhead[RE]{\textit{Harald Høffding: New Investigations into the Primitive Family}}
\fancyhead[LO]{\textit{Tilskueren, vol.\ 5 (1888)}}

\setstretch{1.3}

\title{\textbf{New Investigations into the Primitive Family}\\[8pt]
  \large [Nye Undersøgelser om den primitive Familje]}
\author{Harald Høffding}
\date{\textit{Tilskueren. Maanedsskrift}, vol.~5 (1888), pp.~787--795\\[10pt]
  \small Translated from the Danish}

% ---- original-page markers ----
% \opage{N} marks where printed page N begins, at the same joint as in
% transcription.tex (mid-sentence where the page turns mid-sentence).
\usepackage{marginnote}
\newcommand{\opage}[1]{\marginnote{\footnotesize\textit{[#1]}}}
\newcommand{\apage}[1]{\marginnote{\footnotesize\textit{[#1]}}}

\begin{document}
\maketitle
\thispagestyle{empty}
\clearpage

% --- p. 787 ---
\opage{787}\begin{center}\textbf{\large New Investigations into the Primitive Family.}\end{center}

\begin{center}\small C. N. \emph{Starcke}: Die primitive Familie in ihrer Entstehung und Entwickelung. Leipzig. F. A. Brockhaus. 1888. (Internationale wissenschaftliche Bibliothek. 66. Bind).\end{center}

\begin{center}\rule{1.3cm}{0.4pt}\end{center}

It will hardly surprise anyone that a philosophical author should
occupy himself with investigations into the history of the family. If philosophy
is to be something other and more than constructions proceeding from subjective
consciousness, and if ethics in particular is to be something more than
the setting up of postulates, then a factual material must be procured, and for
ethics that means a historical material, one that
can in many ways shed light on the worth of the principles and on the possibilities
of their application. To be sure, a definite distinction must be drawn --- at least
in my view --- between historical and philosophical
ethics. The former can only present ethical life and ethical
institutions as they have \emph{in fact} developed under certain conditions;
but a task different from that is their \emph{valuation},
the assessment of their significance for the continued, conscious development
of human life. Yet conscious intervention
always presupposes possibilities of development and germs that have come into being
spontaneously. Knowledge of the historical forms of life of the past is
therefore a necessary condition for a fruitful treatment of
ethics. And it is not always a given that historians and ethnologists
arrange the material in the way philosophy needs it. That
requires not only historical criticism but also a philosophical eye for the
decisive points. Just as the history of law cannot do without
the collaboration of the circle of jurists, so philosophical training will
be a necessary presupposition for the treatment and presentation
of historical ethics.

In recent years special attention has been aroused by the primitive
family, i.e.\ the family as it appears at the
% --- p. 788 ---
\opage{788}lowest stages of human development known to us. Whereas
one formerly stopped at the patriarchal family as
one knew it from the Old Testament narratives and from
the Indian poems and law books, the significant step has now
been taken that historical or comparative ethics also began
to take account of the family relations of the savage peoples. It was
now seen that one had no right to regard the Hebrew and
Indian stage of development as the most primitive. There came here a
revolution in thinking similar to the one when linguistics went
beyond Sanskrit and took account of the languages of the savage nations,
% TR: p. 788 the print's „Henstillinger“ (requests, recommendations) makes no sense here (as the transcription notes); rendered as „Forestillinger“, "representations"
or when mythology brought the religious representations of the savages within
its purview. One came upon new and surprising facts
and was thereby led to bold hypotheses. Dr. Starcke's
work, which is built on an extraordinarily thorough knowledge of the
available material, and which bears witness to critical ability
and a sound psychological eye, will quite certainly come to occupy
a considerable place in the literature on the interesting subject
it treats. It disdains the outward attractions that for many
might have lain in the piquant features and entertaining anecdotes
which the subject and the sources offer in abundance.
% TR: p. 788 printer's error „nøgterne“ (for „nøgternt“, beside „strængt“); rendered "sober"
It is sober and strict in its form, but attractive and instructive
for anyone who wishes to be oriented in the subject and to have guidance in
taking a position toward the bold, often wild assertions and consequences
% TR: p. 788 „hvortil det har givet Anledning“ -- „det“ taken as the subject (Æmnet), made explicit
to which the subject has given rise. --- Let me try
briefly to indicate what it is about and what position
Dr. Starcke takes.

It has been noticed that among a great many peoples
around the earth kinship is reckoned not through the father
but through the mother, so that a man's nearest heirs
are not his sons but his sister's children. This custom is found in all
parts of the world. It was then natural to assume that there was a
common cause of it. The great looseness and freedom in sexual
matters that is likewise found among most savage peoples
led people to think of the old maxim: pater
incertus, mater semper certa. (One may be uncertain who
the father is, but the mother is always certain). It was assumed that the
original condition in sexual matters had been a complete
unrestraint and a constant changing, an indiscriminate mating
(promiscuity), which made it impossible to assign the children to
particular fathers and compelled one to take the female line as the basis.
% --- p. 789 ---
\opage{789}The first to put forward the hypothesis that there had been a time
when kinship was reckoned through the woman, an arrangement that marked
an advance by replacing or at any rate regulating promiscuity,
% TR: p. 789 Bachofen's title kept exactly as printed, including the printer's error „Untermehung“ (for Untersuchung) and the Danish ø in „religiøsen“ (for religiösen)
was \emph{Bachofen}, in his work ``Das Mutterrecht. Eine
Untermehung über die Gynaikokratie der alten Welt nach ihrer
religiøsen und rechtlichen Natur'' (Stuttgart 1861). He connected
the predominance of the female line in a certain period with
a bold mythological theory and held that in this period woman
had been highly esteemed, and that her dominion (the matriarchate)
had had great significance for the history of culture by softening the mores
and restraining the animal instincts. He believed he found a support for his theory
in polyandry, the form of sexual cohabitation in which
one woman is the wife of several men. He regarded it as a necessary
intermediate form through which the human race has worked its way
% TR: p. 789 printer's error „Promiskniteten“ (n for u); rendered "promiscuity"
out of promiscuity and approached monogamy. --- Among
those who have set up similar theories, \emph{Mc Lennan} and
\emph{Lubbock} must be named. These investigators proceed with greater criticism
and less mysticism than Bachofen, and Mc Lennan in particular seeks
acutely to show how a good many mores and customs
at higher stages of development, which would otherwise be inexplicable, can
be understood as remnants (survivals) of the female line or of polyandry with
promiscuity as background. Such customs are, e.g., the levirate,
the custom that a man was obliged to marry his deceased
brother's widow, in such a way that the children who sprang from this marriage
were nevertheless reckoned as the deceased's children; and Niyoga, the custom
that the husband (still among the Indians, the Spartans, indeed among the
Germanic peoples of the Middle Ages) could leave his wife to other men and yet
regarded the children thus begotten as his own.

On this theory the picture of the development of the human race,
as regards the relation between man and woman, would to be
sure round itself off in a clear and simple way. From an original
unbound community (promiscuity) we would come to the legalized
community (the female line and polyandry), and from there through
the patriarchal family to monogamous marriage. This
clarity and simplicity has probably exercised a seductive power over
the investigators named. A development may well have taken place
without its course being so easy to survey. And first and foremost
the question remains whether the stages named can really be shown
to be historically given forms of development, each of which has
been dominant in its time.

% --- p. 790 ---
\opage{790}Dr. Starcke denies both that one can infer back from the female line
to an original promiscuity, and that the female line
itself has been as general as the theory in question asserts.
He then seeks to show that the female line, where it prevails,
may have arisen from different motives among the different peoples,
and that likewise the other mores and customs that have been appealed to
must be explained otherwise than according to the theory.

As regards the first point (the female line as proof of
original promiscuity), the author might have mentioned that there are
investigators who assume the generality of the female line at a certain
stage and yet reject the assumption of a general promiscuity.
Thus \emph{Post} in his ``Afrikanische Jurisprudenz. Ethnologisch-juridische
Beiträge'' (1887). This alone suggests that the two
assumptions do not stand and fall together. But most important is the
general consideration that Dr. Starcke puts forward and substantiates:
that from the way in which kinship is reckoned no
conclusion can be drawn as to the nature of marriage. The female line shows
only that regard for the father is set aside, but the question
of the reason for this still stands open. That it cannot (everywhere at
least) have been that it was uncertain who was
the father can be seen from the fact that the female line still to this day holds
under conditions in which the father can very well be identified. The real
reason will be discovered when we consider that the family at primitive
stages does not exist as an independent society, but only as part
of the clan and of the tribe. The group to which the child at
birth attaches and to which it is reckoned is the clan. The child had to be reckoned
\emph{either} to the father's \emph{or} to the mother's clan, and for various
reasons among many peoples it was reckoned to the mother's. But
the family relation proper was not necessarily affected thereby.
The father could stand in the tenderest relation to his children even if
they were reckoned to the mother's clan, not to his, and the same
held for the mother where the male line prevailed.
The determination of kinship is thus a determination of clan, in which
at primitive stages no regard is taken to the relation of procreation.
If we assume that the man, by virtue of his physical strength, was originally
the dominant one, and that women and children were regarded as
his property, the development of the female line can be explained by the
different conditions operating among different peoples. Among the
American peoples (e.g.\ the Columbians) a connection can be traced
% TR: p. 790/791 the page turns inside „Udform-|ning“ (formation); the English turns before "formation"
between the development of the female line and the firmer
% --- p. 791 ---
\opage{791}formation of the clan. To the woman's clan are also reckoned the children she bears,
which from the first are so closely bound up with her;
she is the property of her kin, and her children in turn cannot
be separated from her. The woman's value as a slave, or at any rate as a help
% TR: p. 791 „maa her vel fastholdes“ -- read with the following clause (the clan would not give up that value): the value is retained by the clan; „maa ... vel“ = must presumably (independent review)
in work, must presumably be retained here; the clan would not
renounce it because she entered into sexual union with a man
of another clan. With men of her own clan she was
forbidden, on account of the prevailing exogamy, to enter into
union. Among the African peoples (e.g.\ the Bechuanas)
the motive for the rise of the female line is to be sought in the prevailing
polygamy in connection with the arrangement that each wife lives by
herself, and that the children thus come to form a household
together with the mother, in contrast to the other households, which
together with them gathered around the common head of the patriarchal family.
The position the child obtains in the patriarchal family will here
depend essentially on the mother's rank and dignity. The sons
bear their mothers' names as long as the father lives, and only
after the father's death does the eldest son of the foremost wife enter
into his position and name. The germ of the female line lies in the weight
the mother's position has acquired. --- For the rest, other conditions too will
be of decisive significance. An agricultural people will have
a tendency toward the female line, because the daughter there is of value as a
worker; a cattle-raising people will have a tendency toward the male line,
because there it will be advantageous to sell the daughter and get cattle
instead. For at primitive stages woman is, after all, regarded essentially as
a piece of property that is to be employed as advantageously as possible,
and this view is followed both on the father's and on the husband's
side.

Modern observers and investigators are --- according to Dr. Starcke
--- too inclined to transfer their own feelings to the savages. For
us the consciousness of being the child's physical origin, of its
being the father's blood that flows in the child's veins, is a necessary
component of the paternal feeling proper. But it is not so
for the savage. His representations connect more by merely
external connection than according to real causality. Visual representations
are the dominant ones in all of us and are altogether determining for the
savage, so that what he \emph{sees together} (and only that) also for
him \emph{belongs together}. The clan lives together in its own part of
the village, and the children live with their mother in her own
hut. This spatial being-together becomes of great significance for
% --- p. 792 ---
\opage{792}the development of the female line. And it also explains the remarkable
phenomenon that the savage and the barbarian regards himself as father
even where he \emph{knows} that he is not the physical father. His
wife's children are always his own children as well. Polyandry, levirate
and Niyoga find their explanation easily and naturally in this way of thinking,
without any need to postulate an original indiscriminate
mating. Among the savages the paternal relation is predominantly a legal relation,
not so much a blood relation.

Besides this important psychological factor there is yet
another consideration which according to Dr. Starcke is of decisive significance
for a correct conception of the family institutions of the savages. The
looseness and unrestraint in sexual matters that is found among the
savages has naturally struck observers and investigators, but
has had too great an influence on their theories. Satisfaction
of the sexual instinct plays a large role in the savage's life, but is
also easy to attain, since no narrow bounds are set to it
by the prevailing mores. But precisely for that reason the sexual
relation to another individual does not, among the savages, become of great practical
significance, does not become the basis of an institution. The family as
an institution has therefore from the first and for long ages been ordered
according to quite other considerations than the natural impulses themselves that bring
man and woman together. Only gradually, as the life of representation
develops and it is no longer merely the crude spatial representation
that dominates consciousness, will the paternal relation be regarded essentially
as a blood relation, and with that the demand of chastity on the woman
within, and soon also before, marriage will be given; from the woman
it will then also be extended to the man.

Dr. Starcke ought perhaps to have carried the interesting psychological
consideration indicated here somewhat further. With
higher spiritual development and the finer sense that follows from it for
the unity and coherence of personal life, it will also become
impossible to conceive the sexual instinct and its satisfaction as something
that stands without any connection with the whole of the rest of the personality's
life. Through many fine threads this side of life hangs
together with life's other sides. This shows itself especially in the fact
that the satisfaction of the physical instinct is not the exclusive
goal, but that this satisfaction stands as a good only when
a deeper sympathy with another individual, joy and delight in
% TR: p. 792 printer's error „Kønsinktinktet“ (for „Kønsinstinktet“); rendered "the sexual instinct"
% TR: p. 792/793 the page turns inside „ud-|vikler sig“ (develops); the English turns before "develops"
that individual, thereby receives its seal and its expression. The sexual instinct then
% --- p. 793 ---
\opage{793}develops into the feeling of erotic love proper, and the physical union
into being a spiritual union as well.

People have often read too much romance into the savages
and found in their less restrained sexual life the freedom which, it is
thought, the feeling of erotic love lacks in civilized life. Dr. Starcke,
on the other hand, perhaps reads too little romance into the savages.
In his book the relation between man and woman at
the primitive stages of development appears partly as a merely physical sexual relation,
partly as a purely juridical and economic relation. It would have
been of great interest if he had used his large material
to investigate how far the feeling of erotic love can be demonstrated among
primitive human beings, and what forms it takes. It would
certainly have been shown thereby that the conditions are after all not quite so
different as one might gather from Dr. Starcke's presentation.
Just as it is a great misunderstanding to believe that
the savages live a free and unbound life, since they are even greater
slaves of traditions, ceremonies and rules than most civilized
peoples, so it would also be a misunderstanding to
assume that even the lowest human stage of life we now know
presents no emotional phenomena at all akin to
what we now (as distinct from the mere sexual instinct) call the feeling of erotic love.
But it would certainly be shown that the connection between
feeling and instinct is still closer among the savages than it needs
to be where a higher spiritual development has taken place. The
savage knows no lyric except as an incitement to action or
as a means of reaching his goal. An author who has occupied himself
with \emph{one} significant side of the savages' spiritual life says:
% TR: p. 793 Andrew Lang is quoted in Danish; translated from Høffding's Danish, the English original of Myth, Ritual and Religion not reconstructed
``If you ask an Indian for a love song, he will answer
that a love potion works better. The savage is extraordinarily practical.
His arts, mime and drawing, do not exist for art's own
sake, but as means to a definite end, as methods by
which he can attain something he desires. The young savage whom
Kohl [a traveler in North America] knew put his trust in
having a picture of himself and a picture of his beloved. The
heart of the female picture he smeared with magic powder; this
was, he said, a common custom, and to it were added elegiac
or malicious songs and downright criminal incantations.''
(\emph{Andrew Lang}: Myth, Ritual and Religion. London 1887).

The glaring contrast which the life of the savages nevertheless undeniably
% TR: p. 793/794 the page turns inside „i mere eller mindre flygtig Opblussen | fremtrædende“; the English turns after "fleeting"
displays between the sexual instinct, appearing in more or less fleeting
% --- p. 794 ---
\opage{794}flare-ups, and the fixed, lasting forms for the life and development of the race,
can only be overcome and reconciled through the healthy development of the
feeling of erotic love. The wonderful thing about this
feeling is its power to bind so many of life's forces
around one center, to unite the most ideal and the most elementary
in human nature, and at the same time its power to undergo
metamorphoses through changing fortunes and periods of life. Without
this richness and this power of transformation it would not deserve
the eminent place that poetry has given it.

We have not yet come far beyond the standpoint of the savages
as regards this side of life. This shows itself already
in the fact that so many --- both of the adherents of ``free love''
and of its opponents --- at the word ``marriage'' think first and
foremost of a union between man and woman entered into in legal
or ecclesiastical form. This has turned much bitter
strife into a battle over mere words and has diverted attention
from the real main point, namely the possibility of faithfulness and
inwardness and their significance in the relation between
man and woman. ---

In Dr. Starcke's view there is a great and decisive
difference between the \emph{position} of the family at primitive and at higher
stages of development. He shows in a striking way that the family proper
(the group of husband, wife and children) only at a higher
civilization becomes an independent society, independent of clans and tribes.
Among the savages the family exists only as part
of the clan. This had to be so on account of the unceasing
struggle for existence. The clan persists because association
is necessary in this struggle; in the family, on the other hand,
the goods fought for and acquired are enjoyed. The more, therefore, the struggle for
physical existence recedes, the more will the clan lose in
significance, and the family develop independently. In family life
man comes to partake of an enjoyment of a different kind from that which
arises from the satisfaction of physical instincts, namely ``the
moral enjoyment that life for and with others affords''. --- But
as regards the \emph{motives} that found marriage, there is, according to Dr.
Starcke, not, and (he thinks) there ought not to be, any decisive
difference between primitive and properly human life. ``We
have seen,'' he says toward the end of his book, ``that it is originally
no tender feeling, at least not love,
that instills in the man the wish to marry, that primitive
% --- p. 795 ---
\opage{795}marriage, hard and dark like primitive life, springs from
the most concrete and sober necessities. And still to this day
that marriage is bad in which the erotic is the prominent element
in the relation between the spouses. The common home, in which
each of them has his task, and the common interest in having
children and bringing them up, these were the two foundations on which
marriage was originally built. And from the sympathy that
inevitably springs from the common interests there grows in turn the
love that finally completes and consolidates marriage.'' But
will not precisely the feeling of erotic love more and more become a necessary
presupposition for a man and a woman to be minded
to build on those two foundations? And can it be dispensed with as
a cement that holds together the building that is raised? That it,
in order to work in this way, must undergo a whole process of development
in relation to the given tasks and periods of life has already been
suggested. ---

Although I have missed in Dr. Starcke's book something that in my eyes is significant,
% TR: p. 795 printer's error: a full stop after „Bog“ before lower-case „skilles“ (for a comma); rendered as one sentence
I nonetheless part from it with full recognition
of the thoroughness and acuteness to which it bears witness. Only
imperfectly have I been able in this review to suggest the
wealth of individual investigations that the book contains. It has
appeared in German because so specialized an investigation could not
count on sufficiently many readers in our mother tongue.
But this circumstance ought not to exclude the author from due
recognition from his countrymen.

\begin{flushright}\textsc{Harald Høffding.}\end{flushright}
\begin{center}\rule{2.6cm}{0.4pt}\\[-10pt]\rule{2.6cm}{0.4pt}\end{center}

\end{document}
